% =========================================================
% 8. read() ORM Method
% =========================================================
\section{\texttt{read()} ORM Method}

\begin{frame}[standout]
  \texttt{read()} ORM Method
\end{frame}

\begin{frame}[fragile]{\texttt{READ()} METHOD}
  \begin{itemize}
    \item The \texttt{read()} method is an ORM method used to \textbf{fetch field values}
          from a recordset as plain Python dictionaries.
    \item It is commonly used by the web client and RPC calls.
    \item The method returns a \textbf{list of dictionaries}, not a recordset.
  \end{itemize}
  \begin{block}{Method Syntax}
    The standard syntax of the \texttt{read()} method is:
\begin{lstlisting}
def read(self, fields=None, load='_classic_read'):
    return super().read(fields=fields, load=load)
\end{lstlisting}
    Important parts: \texttt{self}, \texttt{fields}, \texttt{load}.
  \end{block}
\end{frame}

\begin{frame}[fragile]{BREAKING DOWN THE METHOD SIGNATURE}
  \begin{itemize}
    \item \texttt{self}: The recordset whose values you want to read.
    \item \texttt{fields}: An optional list of field names to fetch.
  \end{itemize}
\begin{lstlisting}
fields = ['name', 'age', 'email']
\end{lstlisting}
  \begin{itemize}
    \item If \texttt{fields} is omitted, Odoo reads all readable fields.
    \item Reading only needed fields is faster and cleaner.
    \item \texttt{load}: Controls how relational values are returned
          (classic read vs plain IDs).
  \end{itemize}
\end{frame}

\begin{frame}[fragile]{WHAT DOES \texttt{READ()} RETURN?}
  \begin{itemize}
    \item \texttt{read()} returns a list of dictionaries.
  \end{itemize}
\begin{lstlisting}
student = self.env['student.student'].browse(15)
data = student.read(['name', 'age'])
print(data)
\end{lstlisting}
  Output is (example):
\begin{lstlisting}
[{'id': 15, 'name': 'Ahmed', 'age': 22}]
\end{lstlisting}
  \begin{itemize}
    \item Each dictionary represents one record.
    \item The key \texttt{id} is always included.
  \end{itemize}
\end{frame}

\begin{frame}[fragile]{READING MANY RECORDS}
\begin{lstlisting}
students = self.env['student.student'].browse([15, 16])
data = students.read(['name', 'email'])
\end{lstlisting}
  Output is (example):
\begin{lstlisting}
[
    {'id': 15, 'name': 'Ahmed', 'email': 'ahmed@example.com'},
    {'id': 16, 'name': 'Amina', 'email': 'amina@example.com'},
]
\end{lstlisting}
  \begin{itemize}
    \item Useful when you need serializable data (JSON / RPC / API).
    \item For normal business logic inside Odoo, prefer recordset access:
          \texttt{student.name}.
  \end{itemize}
\end{frame}

\begin{frame}[fragile]{\texttt{READ()} VS RECORDSET ACCESS}
  \begin{enumerate}
    \item \textbf{Recordset access} (preferred in Python business code):
\begin{lstlisting}
name = student.name
age = student.age
\end{lstlisting}
    \item \textbf{read()} (preferred for RPC / exporting dict data):
\begin{lstlisting}
data = student.read(['name', 'age'])[0]
\end{lstlisting}
  \end{enumerate}
  \begin{itemize}
    \item Recordset access keeps ORM features (prefetch, cache, relational browsing).
    \item \texttt{read()} gives simple data structures.
  \end{itemize}
\end{frame}

\begin{frame}[fragile]{METHOD CALL}
  \textbf{Method Call}
\begin{lstlisting}
data = self.env['student.student'].browse(15).read([
    'name',
    'birth_date',
    'gender',
])
\end{lstlisting}
  \begin{itemize}
    \item Here, \texttt{.read(...)} is the method call.
    \item Return type reminder: \textbf{list of dicts}, not a recordset.
  \end{itemize}
\end{frame}
